\documentclass[11pt,a4paper]{article}
\usepackage[margin=1in]{geometry}
\usepackage[T1]{fontenc}
\usepackage[utf8]{inputenc}
\usepackage{lmodern}
\usepackage{parskip}

\begin{document}

\begin{center}
{\Large \textbf{Event Benefits Statement}}\\
{\large ANRF International Travel Support Application}\\[0.5em]
\end{center}

\textbf{Applicant:} [Applicant name hidden], PhD Research Scholar, Department of Chemical Engineering, Indian Institute of Technology Delhi

\textbf{Event:} CCP5 Summer School 2026, Newcastle upon Tyne, United Kingdom, 19 July--30 July 2026

\textbf{Poster Title:} Hydrogen Storage in Pillared Graphene-CNT Frameworks: Junction Topology versus Interlayer Accessibility

CCP5 Summer School 2026 was directly relevant to my doctoral research on molecular dynamics simulation of carbonaceous materials for hydrogen-storage applications. My PhD work focuses on atomistic simulation of graphene, carbon nanotubes, and graphene-CNT frameworks, with the objective of understanding how nanoscale structure controls hydrogen adsorption and storage performance.

My accepted poster presented molecular dynamics simulations of pillared graphene-CNT frameworks at 77 K and 10 bar. The work compared pristine graphene, junction-bearing monolayers, and pillared bilayer frameworks to separate the effects of local graphene-CNT junction chemistry from larger-scale interlayer accessibility. The results identified accessible interlayer volume as a dominant design factor for improving hydrogen uptake in 3D graphene-CNT hydrogen-storage frameworks.

In addition to the poster session, I planned to attend lectures, hands-on computational sessions, and advanced training modules at the summer school. These sessions were directly useful for improving my simulation workflows, especially molecular simulation methodology, Python-based analysis, reproducible computational workflows, and modern approaches used in atomistic materials modelling.

Participation in the event was expected to provide three concrete benefits:

\begin{enumerate}
  \item \textbf{Research feedback:} Presenting the poster would allow me to receive feedback from an international molecular-simulation community on the modelling assumptions, adsorption analysis, and interpretation of hydrogen-storage trends in graphene-CNT frameworks.
  \item \textbf{Training benefit:} The lectures and practical sessions would strengthen my command of simulation methods and scientific computing workflows directly used in my PhD research.
  \item \textbf{Collaboration and future work:} Discussions with researchers working on molecular simulation and computational materials science would help refine my ongoing PhD work and identify future directions for carbon-based hydrogen-storage materials.
\end{enumerate}

The knowledge, feedback, and scientific interactions gained from the summer school would contribute directly to my doctoral thesis, submitted manuscript, and future computational materials research on hydrogen storage in carbon allotropes.

\vfill

\noindent [Applicant name hidden]\\
PhD Research Scholar\\
Department of Chemical Engineering\\
Indian Institute of Technology Delhi

\end{document}
